\chapter*{Abstract}

\vspace{1cm}

This thesis investigates multitask learning for forest monitoring from Sentinel-1 SAR imagery. The main question is whether a single model can estimate multiple forest attributes more efficiently than separate single-task models, and how physics-aware constraints can improve learning.Two datasets were used for experiments. First, the TreeSatAI benchmark was extended with canopy height labels from German LiDAR data, creating TreeSatAI-CHM for joint species classification and height regression. Second, the iMAESTRO dataset provided genus segmentation, height, and biomass labels, enabling tests of allometric constraints that couple height and biomass predictions. The experiments show that multitask learning achieves comparable or slightly better performance than single-task baselines while using fewer parameters and less inference time. Gradient analysis reveals that tasks generally align in deeper network layers while occasionally conflicting in early layers. Physics-aware constraints successfully guide related tasks (height and biomass) toward consistent predictions that respect ecological relationships. The main challenges are balancing tasks with different loss scales and managing task asymmetry. Uncertainty weighting and gradient analysis both address these issues effectively. While absolute performance improvements are modest due to C-band SAR's inherent limits, multitask learning offers a practical solution for operational forest monitoring where computational efficiency matters.
